%% ============================================================
%%  FABRIC — Análisis y Diseño · Versión 6.0
%%  fabricsoft.com.mx · Mayo 2026
%%  Equipo A — Edi + Gerardo · Supervisión Julio Álvarez
%%  v6.0: estado post-demo, correcciones de stack y avance real
%% ============================================================

\documentclass[11pt, a4paper]{article}

%% — Paquetes fundamentales
\usepackage[utf8]{inputenc}
\usepackage[T1]{fontenc}
\usepackage[spanish, es-tabla]{babel}
\usepackage{geometry}
\usepackage{microtype}
\usepackage{parskip}
\usepackage{setspace}

%% — Tipografía
\usepackage{lmodern}
\usepackage{inconsolata}          % Para bloques de código / mono

%% — Color y gráficos
\usepackage[table, dvipsnames]{xcolor}
\usepackage{graphicx}
\usepackage{tikz}

%% — Tablas
\usepackage{booktabs}
\usepackage{longtable}
\usepackage{array}
\usepackage{multirow}
\usepackage{colortbl}
\usepackage{tabularx}

%% — Código
\usepackage{listings}
\usepackage{fancyvrb}

%% — Listas y enumeraciones
\usepackage{enumitem}

%% — Hipervínculos
\usepackage[hidelinks, colorlinks=true, linkcolor=accent, urlcolor=accent]{hyperref}

%% — Encabezados y pies
\usepackage{fancyhdr}
\usepackage{lastpage}

%% — Cajas y frames
\usepackage{mdframed}
\usepackage{tcolorbox}
\tcbuselibrary{skins, breakable}

%% — Títulos y secciones
\usepackage{titlesec}
\usepackage{titletoc}

%% ============================================================
%%  PALETA DE COLORES — tokens del sistema de diseño FABRIC
%% ============================================================

\definecolor{bgbase}     {HTML}{050505}   % Fondo base
\definecolor{bgpanel}    {HTML}{131313}   % Cards y paneles
\definecolor{bgelevated} {HTML}{1A1A1A}   % Elementos flotantes
\definecolor{bordercol}  {HTML}{2A2A2A}   % Bordes secundarios
\definecolor{borderstr}  {HTML}{353535}   % Bordes activos / hover
\definecolor{textpri}    {HTML}{F5F5F5}   % Texto principal
\definecolor{textsec}    {HTML}{8A8A8A}   % Metadatos y labels
\definecolor{textter}    {HTML}{5A5A5A}   % Placeholders
\definecolor{accent}     {HTML}{C9A96E}   % Champagne — único acento
\definecolor{accent2}    {HTML}{A07845}   % Bronce — hover
\definecolor{danger}     {HTML}{B85450}   % Errores y alertas

%% ============================================================
%%  GEOMETRÍA
%% ============================================================

\geometry{
  left   = 2.8cm,
  right  = 2.8cm,
  top    = 3.0cm,
  bottom = 3.0cm,
  headheight = 14pt
}

%% ============================================================
%%  ESTILOS DE SECCIÓN
%% ============================================================

\titleformat{\section}
  {\large\bfseries\color{accent}}
  {§\thesection}{1em}{}[\vspace{2pt}\hrule]

\titleformat{\subsection}
  {\normalsize\bfseries\color{textpri}}
  {\thesubsection}{1em}{}

\titleformat{\subsubsection}
  {\small\bfseries\color{textsec}}
  {\thesubsubsection}{1em}{}

%% ============================================================
%%  ESTILOS DE LISTADO (CÓDIGO)
%% ============================================================

\lstset{
  basicstyle   = \small\ttfamily\color{textpri},
  backgroundcolor = \color{bgpanel},
  frame        = single,
  rulecolor    = \color{bordercol},
  keywordstyle = \color{accent},
  stringstyle  = \color{accent2},
  commentstyle = \color{textsec}\itshape,
  breaklines   = true,
  showstringspaces = false,
  tabsize      = 2
}

%% ============================================================
%%  CAJAS PERSONALIZADAS
%% ============================================================

%% Caja de doctrina / principio
\newtcolorbox{doctrinabox}[1][]{
  colback   = bgpanel,
  colframe  = accent,
  coltitle  = bgbase,
  fonttitle = \small\bfseries\ttfamily,
  title     = {#1},
  left      = 8pt,
  right     = 8pt,
  top       = 6pt,
  bottom    = 6pt,
  breakable
}

%% Caja de aviso / pendiente
\newtcolorbox{pendiente}{
  colback  = bgpanel,
  colframe = danger,
  left     = 8pt,
  right    = 8pt,
  top      = 4pt,
  bottom   = 4pt,
  breakable
}

%% Caja de implementado
\newtcolorbox{implementado}{
  colback  = bgpanel,
  colframe = accent,
  left     = 8pt,
  right    = 8pt,
  top      = 4pt,
  bottom   = 4pt,
  breakable
}

%% Caja de nota técnica
\newtcolorbox{notatec}[1][Nota técnica]{
  colback   = bgpanel,
  colframe  = bordercol,
  coltitle  = accent,
  fonttitle = \small\bfseries\ttfamily,
  title     = {#1},
  left      = 8pt,
  right     = 8pt,
  top       = 4pt,
  bottom    = 4pt,
  breakable
}

%% ============================================================
%%  ENCABEZADO Y PIE DE PÁGINA
%% ============================================================

\pagestyle{fancy}
\fancyhf{}
\renewcommand{\headrulewidth}{0.4pt}
\renewcommand{\headrule}{\hbox to\headwidth{\color{bordercol}\leaders\hrule height \headrulewidth\hfill}}
\fancyhead[L]{\small\ttfamily\color{textsec} FABRIC · Análisis y Diseño V6.0}
\fancyhead[R]{\small\ttfamily\color{textsec} fabricsoft.com.mx}
\fancyfoot[C]{\small\ttfamily\color{textsec} \thepage\ / \pageref{LastPage}}

%% ============================================================
%%  COMANDOS DE UTILIDAD
%% ============================================================

\newcommand{\token}[1]{\texttt{\color{accent}#1}}
\newcommand{\ruta}[1]{\texttt{\color{accent2}#1}}
\newcommand{\estado}[1]{\textbf{\color{accent}#1}}
\newcommand{\pendientetag}{\textbf{\color{danger}[PENDIENTE]}}
\newcommand{\implementadotag}{\textbf{\color{accent}[IMPLEMENTADO]}}
\newcommand{\parcialtag}{\textbf{\color{accent2}[PARCIAL]}}

%% ============================================================
%%  DOCUMENTO
%% ============================================================

\begin{document}

%% ——— PORTADA ———————————————————————————————————————————

\begin{titlepage}
\pagecolor{bgbase}
\color{textpri}

\vspace*{3cm}

\begin{center}
  {\Huge\bfseries FABRIC}\\[0.4em]
  {\large\ttfamily\color{accent} ORACLE CRITICAL ENGINEERING}\\[0.2em]
  {\small\ttfamily\color{textsec} fabricsoft.com.mx}
\end{center}

\vspace{2cm}

\begin{center}
  \begin{tikzpicture}
    \draw[color=accent, line width=0.8pt] (0,0) -- (\textwidth,0);
  \end{tikzpicture}
\end{center}

\vspace{1.5cm}

\begin{center}
  {\LARGE\bfseries ANÁLISIS Y DISEÑO}\\[0.6em]
  {\Large VERSIÓN 6.0}\\[0.4em]
  {\normalsize\color{textsec} Propuesta técnica consolidada · Estado post-demo}
\end{center}

\vspace{2cm}

\begin{center}
  \begin{tabular}{ll}
    \texttt{\color{textsec}Equipo}        & Edi + Gerardo\\[0.3em]
    \texttt{\color{textsec}Supervisión}   & Julio Álvarez\\[0.3em]
    \texttt{\color{textsec}Standup}       & 11:00 AM diario\\[0.3em]
    \texttt{\color{textsec}Demo intermedia} & Jueves 21 mayo 2026 {\color{accent}✓}\\[0.3em]
    \texttt{\color{textsec}Demo final}    & Lunes 25 mayo 2026 {\color{accent}✓}\\[0.3em]
    \texttt{\color{textsec}Versión}       & {\color{accent}V6.0 — Correcciones post-demo}\\[0.3em]
    \texttt{\color{textsec}Fecha}         & Mayo 2026\\
  \end{tabular}
\end{center}

\vspace{3cm}

\begin{center}
  \begin{tikzpicture}
    \draw[color=bordercol, line width=0.4pt] (0,0) -- (\textwidth,0);
  \end{tikzpicture}
\end{center}

\vspace{0.8cm}

\begin{center}
  {\small\color{textsec}
  Este documento recoge la propuesta técnica completa del Equipo A.\\
  La v6.0 corrige las desviaciones de stack detectadas entre el documento\\
  original y la implementación real aprobada tras las demos del 21 y 25 de mayo.
  }
\end{center}

\end{titlepage}

\pagecolor{white}
\color{black}

%% ——— NOTA DE VERSIÓN ————————————————————————————————

\begin{notatec}[Cambios respecto a V5.0]
La versión 5.0 describía un stack basado en \textbf{Next.js 15 + App Router}.
Tras las demos y la integración con el equipo, el proyecto se ejecuta sobre
\textbf{Vite 6 + React 19 + React Router v7}. Las secciones afectadas son
§6 (arquitectura frontend), §8 (admin), §9 (integraciones) y §10 (infraestructura).
El posicionamiento, la doctrina, el sistema de diseño y los flujos de negocio
permanecen sin cambios.
\end{notatec}

\vspace{1em}

%% ——— TABLA DE CONTENIDO ——————————————————————————————

\tableofcontents
\newpage

%% ——————————————————————————————————————————————————————
%%  §1  CONTEXTO DEL NEGOCIO Y POSICIONAMIENTO
%% ——————————————————————————————————————————————————————

\section{Contexto del negocio y posicionamiento estratégico}

FABRIC opera en un segmento muy específico: implementaciones Oracle críticas en empresas
que no pueden permitirse incertidumbre operativa. El posicionamiento parte de una premisa
simple pero disruptiva — la industria celebra el go-live como hito final, y FABRIC lo
considera el inicio del trabajo real.

\subsection{Misión}

Acompañar y liderar la remediación de proyectos Oracle críticos garantizando estabilidad
operativa durante el primer ciclo transaccional en producción, erradicando el enfoque
transitorio del go-live.

\subsection{Posicionamiento maestro}

\begin{table}[h]
\centering
\begin{tabularx}{\linewidth}{>{\ttfamily\color{textsec}\small}l X}
\toprule
Elemento & Definición \\
\midrule
Categoría madre   & Oracle Critical Engineering \\
Promesa contractual & \textit{``No entregamos en go-live. Entregamos cuando tu primer ciclo crítico opera en producción.''} \\
Cliente objetivo  & CFO + CTO de empresa USD 50M–500M+ en Servicios Financieros, Inmobiliario o Logística \\
Enemigo narrativo & La implementación Oracle que termina en go-live y abandona al cliente sin operar \\
Tagline           & Oracle Cloud, entregado en producción. \\
Marca pública     & FABRIC — única forma válida \\
Razón social      & FABRIC SOFT MEXICO SA DE CV — solo en footer/legal \\
\bottomrule
\end{tabularx}
\end{table}

\subsection{Objetivo general}

Establecer \texttt{fabricsoft.com.mx} como un gatekeeper de autoridad técnica, no como
un canal de ventas convencional. La plataforma está diseñada para que el directivo
demuestre que su empresa califica para trabajar con FABRIC, y no al revés.

\subsection{Objetivos específicos}

\begin{table}[h]
\centering
\begin{tabularx}{\linewidth}{X >{\small}X}
\toprule
Objetivo & Cómo se mide \\
\midrule
Proyectar autoridad técnica desde el primer scroll &
  Paleta oscura · cero azules · Cormorant Garamond en titulares · quietud en el hero \\
Ejecutar filtros de admisión antes del primer contacto humano &
  100\% de leads con dominio corporativo, cargo C-level y revenue declarado \\
Demostrar capacidad Oracle en tiempo real &
  Agente IA responde consultas Oracle · Kill Switch activo \\
Publicar evidencia verificable &
  Caso APE Plazas · transparencia en 3 niveles · 3 papers con gating \\
Alcanzar Lighthouse 95+ &
  LCP $<$ 1.5s · CLS $<$ 0.05 · SEO 100 \\
Panel admin funcional &
  Leads visibles, aprobación/rechazo, métricas editables \\
\bottomrule
\end{tabularx}
\end{table}

\subsection{FABRIC Lifecycle — las 5 fases del modelo de entrega}

\begin{table}[h]
\centering
\begin{tabularx}{\linewidth}{>{\ttfamily\color{accent}\bfseries}l >{\bfseries}l X X}
\toprule
Fase & Nombre & Descripción & Por qué importa \\
\midrule
01 & DIAGNOSE   & Análisis ejecutivo de arquitectura, procesos y riesgos &
                  Sin diagnóstico honesto, cualquier propuesta es una apuesta \\
02 & ARCHITECT  & Diseño técnico con FSO Engine, roadmap 30-60-90 días &
                  La solución se diseña antes de tocar el sistema \\
03 & DEPLOY     & Implementación con Zero-Trust y gobierno semanal &
                  Sin gobierno, los proyectos Oracle se expanden sin control \\
04 & STABILIZE  & \textbf{Acompañamiento durante el primer ciclo crítico} &
                  \textbf{El go-live es un hito intermedio. Esta fase es la promesa.} \\
05 & OPTIMIZE   & Optimización continua con agentes IA propios &
                  El ERP mejora con el tiempo en lugar de degradarse \\
\bottomrule
\end{tabularx}
\end{table}

\subsection{Doctrina FABRIC V4.0 — cinco compromisos contractuales}

\begin{doctrinabox}[01 · Entrega en primer ciclo crítico]
El proyecto no se entrega en el go-live. Se entrega cuando el primer cierre contable,
ciclo operativo o ciclo regulatorio crítico opera en producción con estabilidad
documentada y verificada.
\end{doctrinabox}

\begin{doctrinabox}[02 · Solo seniors. Cero juniors facturables.]
Cada consultor asignado tiene mínimo 8 años de experiencia real en Oracle.
Sin excepciones. Verificable en el CMS.
\end{doctrinabox}

\begin{doctrinabox}[03 · Fixed-Price por fase. Cero sorpresas.]
Presupuestos cerrados por fase. Si el atraso es imputable a FABRIC, las semanas
adicionales no se facturan. El Panel Admin lo controla.
\end{doctrinabox}

\begin{doctrinabox}[04 · Cero reportes manuales post go-live.]
Al cierre del primer ciclo crítico, ningún proceso financiero o ejecutivo debe
correr fuera del ERP. Si queda uno por causa atribuible a FABRIC, se resuelve
sin costo hasta eliminarlo.
\end{doctrinabox}

\begin{doctrinabox}[05 · Transición formal con documentación viva.]
Acta firmada por todos los stakeholders. Incluye: tablero de KPIs verificado,
adopción medida, plan de soporte, y documentación auditable por el cliente
sin dependencia de FABRIC.
\end{doctrinabox}

\newpage

%% ——————————————————————————————————————————————————————
%%  §2  ANÁLISIS DE MERCADO
%% ——————————————————————————————————————————————————————

\section{Análisis de mercado y perfiles de decisión}

\subsection{Prioridades comerciales}

FABRIC tiene tres tipos de prospecto. El orden refleja ciclo de venta,
tasa de conversión y urgencia.

\begin{table}[h]
\centering
\begin{tabularx}{\linewidth}{>{\bfseries\ttfamily\color{accent}}l X >{\small}X}
\toprule
Prioridad & Perfil & Conversión esperada \\
\midrule
1 — Rescate &
  Empresa con Oracle Fusion implementado y problemas activos: cierre contable de semanas,
  reportes manuales corriendo en paralelo, usuarios en Excel, incidencias sin resolver
  desde el go-live. &
  Ciclo 3-6 semanas · 40-50\% · ticket USD 150K-500K \\
2 — Migración &
  Empresa con SAP S/4 HANA, Oracle EBS R12, JD Edwards, PeopleSoft o
  Microsoft Dynamics evaluando mover a Oracle Fusion Cloud &
  Ciclo 6-12 meses · 15-25\% · ticket USD 300K-2M \\
3 — Greenfield &
  Empresa sin ERP empresarial, operando en QuickBooks o Excel, crecimiento
  que exige plataforma moderna &
  Ciclo 4-8 meses · 20-30\% · ticket USD 250K-1.5M \\
\bottomrule
\end{tabularx}
\end{table}

\subsection{Mercado objetivo — tres verticales}

\begin{table}[h]
\centering
\begin{tabularx}{\linewidth}{>{\bfseries}l X l}
\toprule
Vertical & Operación crítica que resuelve FABRIC & Referencias \\
\midrule
Servicios Financieros / Fintech &
  Conciliaciones bancarias, compliance CNBV, cierre contable regulatorio &
  Aplazo, ABC Capital \\
Inmobiliario / Centros Comerciales &
  Multi-entidad, rentas variables por arrendatario, consolidación multi-plaza &
  APE Plazas \\
Logística / Distribución &
  Supply chain integrado, trazabilidad fiscal CFDI 4.0, operación multi-CD &
  ALMEX, AJE, ENEX \\
\bottomrule
\end{tabularx}
\end{table}

\subsection{Buyer personas}

\subsubsection{CFO — Comprador económico}

Director Financiero de empresa USD 50M–500M+. Su métrica principal es la
predictibilidad de costos. Su mayor frustración: consultores que celebran el go-live,
cobran y dejan cierres de 15 días con reportes en Excel.

En el sitio: valida la Doctrina (compromisos 01, 03 y 04), revisa APE Plazas,
usa el TCO Comparator para justificar la inversión.

\subsubsection{CTO / CIO — Comprador técnico}

Responsable de la estabilidad del ERP y las integraciones. Mayor dolor: firmas que
venden seniors en la propuesta y mandan juniors sin capacidad de resolver incidentes.

En el sitio: interroga al Agente Oracle para verificar profundidad técnica, inspecciona
FABRIC OS y descarga los papers de investigación.

\subsection{Caso ancla — APE Plazas}

\begin{table}[h]
\centering
\begin{tabularx}{\linewidth}{X >{\color{accent}\bfseries}l}
\toprule
Hito & Estado \\
\midrule
Go-live en fecha contractual: 06 abril 2026         & Cumplido \\
Configuración multi-entidad en producción           & Operativa \\
Cierre contable de abril 2026 sin incidencias       & Ejecutado el 30 de abril \\
Tablero ejecutivo activo para Dirección y CFO       & Activo \\
Adopción de usuarios clave documentada              & Verificada \\
Transición formal a soporte                         & Acta firmada \\
\bottomrule
\end{tabularx}
\caption{Hitos verificables · Publicación pendiente de autorización formal del cliente}
\end{table}

\newpage

%% ——————————————————————————————————————————————————————
%%  §3  ARQUITECTURA DE LA INFORMACIÓN
%% ——————————————————————————————————————————————————————

\section{Arquitectura de la información y navegación}

\subsection{Home — 19 secciones implementadas}

\begin{longtable}{>{\ttfamily\color{accent}\small}l >{\bfseries}l X >{\small\bfseries}l}
\toprule
Sección & Nombre & Contenido & Estado \\
\midrule
\endhead
S01 & Hero & Manifiesto Oracle Critical Engineering · 2 CTAs · sin imágenes & \color{accent}Implementado \\
S02 & Optimizador OCI & Lead magnet · comparador OCI vs. AWS/GCP/Azure & \color{accent}Implementado \\
S03 & Calculadora TCO & TCO ERP Comparator interactivo · modal con 8 campos & \color{accent}Implementado \\
S04 & TCO Waitlist & Waitlist integrada a la calculadora TCO & \color{accent}Implementado \\
ChatIA & Agente Oracle & Terminal interactiva · 3 escenarios · UI completa & \color{accent2}Parcial — API pendiente \\
S05 & Análisis de Fallas & Diagnóstico Oracle Fusion fallido · patrones & \color{accent}Implementado \\
S06 & Doctrina & Los 5 compromisos contractuales · DoctrinaModal & \color{accent}Implementado \\
S06b & Fixed-Price & Explicación del modelo Fixed-Price por fase & \color{accent}Implementado \\
S07 & Casos Ancla & APE Plazas + Aplazo · tabla de hitos · CTA caso completo & \color{accent}Implementado \\
S07b & Rescue Assessment & Diagnóstico de proyecto Oracle fallido · formulario & \color{accent2}UI lista · sin submit real \\
S08 & Industrias & 3 verticales focales con dolor específico por sector & \color{accent}Implementado \\
S09 & FABRIC OS & 4 capas + catálogo de 6 FSOs & \color{accent2}Parcial — capas no expandibles \\
S10 & Lifecycle & 5 fases · STABILIZE resaltada · timeline horizontal & \color{accent}Implementado \\
S11 & Office Hours & Calendario de slots · criterios de admisión & \color{accent2}UI lista · sin Calendly \\
S12 & Referencias & 5 referencias disponibles bajo NDA & \color{accent}Implementado \\
S12b & Criterios & Doctrina de filtrado · admisión y rechazo & \color{accent}Implementado \\
S13 & Transparencia & 3 niveles de métricas verificables & \color{accent}Implementado \\
S14 & Investigación & 3 papers con UI de descarga & \color{accent2}Sin PDFs reales ni gating \\
S15 & Founder + Wait List & Manifiesto · bio Julio · capacidad · admisión & \color{accent}Implementado \\
\bottomrule
\caption{Secciones del home · Mayo 2026}
\end{longtable}

\subsection{Rutas activas}

\begin{table}[h]
\centering
\begin{tabularx}{\linewidth}{>{\ttfamily\color{accent2}}l X >{\small\bfseries}l}
\toprule
Ruta & Descripción & Estado \\
\midrule
/ & Home completa (19 secciones) & \color{accent}Activa \\
/casos/:slug & Detalle de caso por slug (APE Plazas, Aplazo) & \color{accent}Activa \\
/sign-in/* & Autenticación pública Clerk & \color{accent}Activa \\
/sign-up/* & Registro público Clerk & \color{accent}Activa \\
/admin & Dashboard administrativo & \color{accent}Activa · protegida \\
/admin/leads & Gestión de leads & \color{accent}Activa · protegida \\
/admin/metricas & Editor de métricas & \color{accent}Activa · protegida \\
/admin/capacidad & Control de capacidad & \color{accent}Activa · protegida \\
/admin/office-hours & Gestión de slots & \color{accent}Activa · protegida \\
/admin/logs & Audit trail & \color{accent}Activa · protegida \\
/aplicar & Formulario de aplicación (5 campos) & \color{danger}Pendiente \\
/transparencia & Metodología pública de métricas & \color{danger}Pendiente \\
/doctrina/generator & Generador de cláusulas & \color{danger}Pendiente V1.5 \\
/rechazados & Registro de proyectos rechazados & \color{danger}Pendiente \\
/casos/:slug/audit-trail & Timeline público verificable & \color{danger}Pendiente \\
\bottomrule
\end{tabularx}
\end{table}

\subsection{Navegación persistente}

\begin{itemize}[leftmargin=*]
  \item Logo FABRIC en Cormorant Garamond, enlaza al home
  \item Links principales: Doctrina · Casos · Industrias · FABRIC OS · Transparencia
  \item CTA \texttt{Aplicar →}: color champagne \token{\#C9A96E}, subrayado expansivo del centro 300ms
  \item Scroll: navbar transparente en el top, fondo \token{\#131313} al pasar 80px
  \item Mobile desde 968px: hamburger con overlay fullscreen oscuro
\end{itemize}

\newpage

%% ——————————————————————————————————————————————————————
%%  §4  SISTEMA DE DISEÑO
%% ——————————————————————————————————————————————————————

\section{Sistema de diseño y especificaciones UI/UX}

El diseño de FABRIC no es una elección estética. Es una decisión estratégica.
Cada elemento visual comunica selectividad, seriedad técnica y autoridad.

\subsection{Paleta de colores}

\begin{table}[h]
\centering
\begin{tabularx}{\linewidth}{>{\ttfamily\color{accent}\small}l >{\ttfamily}l X}
\toprule
Variable CSS & Valor & Uso \\
\midrule
--bg-base       & \#050505 & Fondo base · negro casi puro \\
--bg-panel      & \#131313 & Cards, terminal, navbar con scroll \\
--bg-elevated   & \#1A1A1A & Elementos flotantes o modales \\
--border        & \#252525 & Separadores y bordes secundarios \\
--border-strong & \#353535 & Bordes en estado activo o hover \\
--text-primary  & \#F5F5F5 & Titulares y cuerpo principal \\
--text-secondary & \#8A8A8A & Metadatos, labels, descripciones \\
--text-tertiary & \#5A5A5A & Placeholders y texto desactivado \\
--accent        & \#C9A96E & Champagne — único acento del sitio \\
--accent-2      & \#A07845 & Bronce — hover del acento \\
--danger        & \#B85450 & Alertas y errores de validación \\
\bottomrule
\end{tabularx}
\end{table}

\begin{notatec}[Corrección v6.0 — color bg-base]
El documento v5.0 especificaba \token{\#0A0A0A} como fondo base. La implementación
aprobada usa \token{\#050505}. El resto de la paleta permanece sin cambios.
\end{notatec}

\subsubsection{Prohibiciones absolutas de color}

Azul saturado tipo SaaS (\texttt{\#3B82F6} y similares) · verde fluorescente · rojo
neón · gradientes visibles · glassmorphism · neumorphism.

\subsection{Tipografía}

\begin{table}[h]
\centering
\begin{tabularx}{\linewidth}{>{\bfseries}l X >{\small}l}
\toprule
Familia & Uso en la interfaz & Rango de tamaño \\
\midrule
Cormorant Garamond & Titulares principales, subtítulos, citas editoriales &
  clamp(48px, 6vw, 88px) \\
Inter & Cuerpo de texto, párrafos, descripciones &
  17px -- 18px \\
JetBrains Mono & Terminal IA, código, tablas de datos, labels, badges &
  11px -- 15px \\
\bottomrule
\end{tabularx}
\end{table}

\begin{notatec}[Corrección v6.0 — tipografía]
El documento v5.0 especificaba \textbf{GT Sectra Display} para titulares.
GT Sectra es una fuente comercial con licencia de pago que requiere contrato
con Klim Type Foundry. La implementación aprobada usa \textbf{Cormorant Garamond}
(licencia SIL OFL, disponible en Google Fonts) como serif principal, manteniendo
el carácter editorial y la elegancia especificada. GT Sectra queda como especificación
ideal para fases futuras con presupuesto de licencias tipográficas.
\end{notatec}

\subsection{Layout y espaciado}

\begin{itemize}[leftmargin=*]
  \item Grid de 12 columnas con gutter de 24px y \texttt{max-width: 1280px}
  \item Espaciado en múltiplos estrictos de 8px: 8 · 16 · 24 · 40 · 64 · 96 · 160px
  \item Padding vertical de secciones: mínimo 160px en desktop (110px en implementación actual con \texttt{.demo-section})
  \item Breakpoint mobile: 968px
\end{itemize}

\subsection{Microinteracciones}

El sitio sigue una política de quietud visual. Las animaciones existen para reforzar
la jerarquía, no para decorar.

\begin{table}[h]
\centering
\begin{tabularx}{\linewidth}{>{\bfseries}l X}
\toprule
Animación & Especificación técnica \\
\midrule
Scroll reveal &
  opacity 0→1 + translateY 24px→0 · 200ms ease-out ·
  stagger 80ms · \texttt{IntersectionObserver} \\
Typewriter en terminal IA &
  30ms por carácter · cursor parpadeante CSS keyframes \\
Subrayado expansivo en CTAs &
  Línea champagne que se expande del centro en 300ms · \texttt{::after scaleX 0→1} \\
Counter del hero &
  0 → valor real al entrar al viewport · 800ms ease-out · solo en primer scroll \\
Navbar al hacer scroll &
  Transparente → \token{\#131313} con \texttt{border-bottom} en 200ms al pasar 80px \\
\midrule
\multicolumn{2}{l}{\textit{Prohibido: parallax · scroll-jacking · confetti · animaciones \textgreater 400ms · más de 2 simultáneas}} \\
\bottomrule
\end{tabularx}
\end{table}

\subsubsection{Regla del Hero}

El brief original especifica \textbf{cero animaciones, quietud total} en el Hero.
La implementación actual incluye globe animado, typewriter y partículas flotantes.
Este punto requiere revisión y aprobación de Julio antes de la versión de producción.

\subsection{Nota de cascada CSS}

\begin{notatec}[Arquitectura CSS — cascada de capas]
El archivo \ruta{maquetado-dossier.css} contiene CSS \textbf{sin capa} (unlayered).
En CSS Cascade Level 5, el CSS unlayered tiene mayor prioridad que cualquier
\texttt{@layer}, independientemente de especificidad. Esto significa que el reset
universal \texttt{* \{ margin: 0; padding: 0 \}} sobreescribía todas las utilidades
de Tailwind (\texttt{px-6}, \texttt{py-24}, etc.) en las secciones de Edilberto.

\textbf{Solución aplicada:} el reset de \texttt{margin/padding} se scopeó a
\texttt{.demo-section *}, y los tamaños de \texttt{h2/h3/h4} se scopearon a
\texttt{.demo-section h2/h3/h4}. Las secciones S07-S15 siguen usando el sistema CSS
del maquetado; las secciones S01-S06 usan Tailwind utilities correctamente.
\end{notatec}

\newpage

%% ——————————————————————————————————————————————————————
%%  §5  MAQUETADO Y PROTOTIPADO
%% ——————————————————————————————————————————————————————

\section{Maquetado y prototipado visual}

\subsection{Entregables producidos}

\begin{itemize}[leftmargin=*]
  \item Home completa en 19 secciones: desktop y mobile
  \item Flujos UI: Rescue Assessment · ERP TCO Comparator · Wait List · Office Hours
  \item Panel admin: dashboard, tabla de leads, editor de métricas, capacidad, logs
  \item Terminal del Agente Oracle: UI completa con 3 escenarios
  \item Archivos de referencia en \ruta{design/} (maquetas HTML y CSS originales)
\end{itemize}

\subsection{Criterios de aprobación visual}

\begin{itemize}[leftmargin=*]
  \item Coherencia con los tokens definidos: fondo oscuro, champagne como único acento,
        bordes sobrios, Cormorant Garamond en titulares
  \item El primer viewport responde las 5 preguntas del tronco sin scroll
  \item Mobile responsive sin scroll horizontal · touch targets mínimos 44px
  \item Hero: tipografía únicamente, sin imágenes (animaciones: pendiente revisión)
\end{itemize}

\newpage

%% ——————————————————————————————————————————————————————
%%  §6  ARQUITECTURA TÉCNICA FRONTEND
%% ——————————————————————————————————————————————————————

\section{Arquitectura técnica — capa de presentación}

\begin{notatec}[Corrección v6.0 — Stack]
La versión 5.0 describía Next.js 15 + App Router. La implementación real
aprobada usa \textbf{Vite 6 + React 19 + React Router v7}. Las implicaciones
técnicas se detallan en esta sección.
\end{notatec}

\subsection{Stack tecnológico}

\begin{table}[h]
\centering
\begin{tabularx}{\linewidth}{>{\bfseries}l >{\ttfamily\small}l X}
\toprule
Herramienta & Versión & Por qué la usamos \\
\midrule
Vite & 6.4.x &
  Dev server ultrarrápido con HMR. Build optimizado con Rollup.
  Configuración simple para SPA \\
React & 19.x &
  Framework UI. Hooks, Context API, componentes funcionales \\
TypeScript & 6.x strict &
  Código auditable, reduce errores en desarrollo \\
React Router & v7.x &
  Enrutamiento client-side. \texttt{<Routes>}, \texttt{<Route>}, \texttt{<Navigate>} \\
Tailwind CSS & 4.3.x &
  Tokens del design system en \ruta{src/index.css} vía \texttt{@theme}.
  Utilities en \texttt{@layer utilities} \\
Framer Motion & 12.x &
  Fade-ins al scroll y typewriter de la terminal. Tree-shakeable \\
Clerk & 5.x &
  Autenticación del panel admin. Sesiones, roles, middleware \\
Sonner & 2.x &
  Toast notifications en flujos de captura \\
@tailwindcss/vite & 4.x &
  Plugin oficial Tailwind para Vite \\
\bottomrule
\end{tabularx}
\end{table}

\subsection{Estructura de carpetas}

\begin{lstlisting}[language=bash, caption={Estructura real del proyecto}]
fabricsoft/                    # Monorepo raiz
|-- src/                       # Frontend Vite / React (app principal)
|   |-- main.tsx               # Entry point
|   |-- App.tsx                # Root component
|   |-- index.css              # Tokens globales + overrides pre-dossier
|   |-- maquetado-dossier.css  # CSS base S07-S15 (unlayered)
|   |-- maquetado-interacciones.css  # CSS de modales
|   |
|   |-- routers/
|   |   `-- AppRouter.tsx      # React Router v7 - rutas publicas y admin
|   |
|   |-- layouts/
|   |   `-- publicLayout.tsx   # Wrapper Header + Footer
|   |
|   |-- pages/
|   |   |-- public/
|   |   |   |-- home/          # 19 secciones: s01-hero ... s15-founder
|   |   |   |-- chat/          # chatIa.tsx - Agente Oracle UI
|   |   |   |-- casos/         # CasoPage.tsx - /casos/:slug
|   |   |   |-- header/        # headerPublic.tsx
|   |   |   `-- footer/        # footerPublic.tsx
|   |   `-- admin/             # Panel administrativo (Clerk)
|   |       |-- AdminDashboard.tsx
|   |       |-- AdminLeads.tsx
|   |       |-- AdminMetricas.tsx
|   |       |-- AdminCapacidad.tsx
|   |       |-- AdminOfficeHours.tsx
|   |       |-- AdminLogs.tsx
|   |       `-- AdminLayout.tsx
|   |
|   |-- components/
|   |   |-- InteractionManager.tsx   # Modales I01-I07
|   |   |-- SectionNavigator.tsx     # Navegacion de secciones
|   |   `-- RetroGrid.tsx
|   |
|   |-- store/
|   |   |-- fabricStore.ts     # Fuente unica de verdad (en memoria)
|   |   `-- FabricContext.tsx  # React Context + hooks
|   |
|   |-- hooks/
|   |   `-- useInViewOnce.ts   # IntersectionObserver para scroll reveal
|   |
|   `-- assets/
|       `-- logo/              # Logo FABRIC
|
|-- frontend/                  # Next.js SSR - Edilberto (paralelo)
|-- backend/                   # API Express + TypeScript (scaffolding)
|-- public/                    # Assets estaticos Vite
|-- docs/                      # Documentacion
`-- design/                    # Maquetas y referencias visuales
\end{lstlisting}

\subsection{Gestión de estado}

El estado global de la aplicación vive en \ruta{src/store/}.
No se usa Redux ni Zustand — React Context API es suficiente para esta fase.

\begin{table}[h]
\centering
\begin{tabularx}{\linewidth}{>{\ttfamily\color{accent2}}l X}
\toprule
Archivo & Responsabilidad \\
\midrule
fabricStore.ts &
  Fuente única de verdad. Define tipos, datos iniciales (en memoria) y
  helpers. Cuando se conecte el backend, \textbf{solo este archivo cambia}. \\
FabricContext.tsx &
  React Provider + hooks: \texttt{useFabric()}, \texttt{useMetrica()},
  \texttt{useCapacidad()}. Métodos: \texttt{cycleSlot}, \texttt{updateMetrica},
  \texttt{updateLeadStatus}, \texttt{reservarSlot}, \texttt{confirmarSlot}. \\
\bottomrule
\end{tabularx}
\end{table}

\begin{pendiente}
\textbf{Persistencia:} actualmente todos los datos (leads, métricas, slots, capacidad)
viven en memoria de sesión. Si se recarga el servidor, los cambios del admin se
pierden. La conexión a MongoDB se implementa en la siguiente fase — solo requiere
modificar \texttt{fabricStore.ts}.
\end{pendiente}

\subsection{Estrategia de render}

A diferencia del stack Next.js descrito en v5.0, la aplicación es una
\textbf{SPA (Single Page Application)} con renderizado client-side.

\begin{table}[h]
\centering
\begin{tabularx}{\linewidth}{>{\ttfamily\color{accent2}}l X}
\toprule
Ruta & Estrategia \\
\midrule
/ (home) &
  Client-side rendering. Todas las secciones se cargan al inicializar la app.
  El counter de proyectos activos viene del store. \\
/casos/:slug &
  Client-side. Datos hardcodeados por slug.
  Migrar a CMS cuando se implemente Sanity. \\
/admin/* &
  Client-side. Protegido por Clerk. Se renderiza solo si la sesión
  está autenticada. \\
\bottomrule
\end{tabularx}
\caption{Sin SSG, ISR ni Server Actions en la implementación actual (Vite SPA)}
\end{table}

\subsection{Diferencias con Next.js (v5.0 vs. v6.0)}

\begin{table}[h]
\centering
\begin{tabularx}{\linewidth}{X >{\color{danger}}X >{\color{accent}}X}
\toprule
Característica & V5.0 (Next.js) & V6.0 (Vite) \\
\midrule
Carga de fuentes & next/font en compilación & Google Fonts vía @import \\
Optimización de imágenes & next/image (WebP automático) & \texttt{<img>} estándar \\
Render & SSG / ISR / SSR por página & SPA client-side \\
Server Actions & Sí & No \\
SEO automático & generateMetadata() & Meta tags manuales \\
Rutas de API & app/api/route.ts & backend/ Express separado \\
\bottomrule
\end{tabularx}
\end{table}

\newpage

%% ——————————————————————————————————————————————————————
%%  §7  MOTOR DE INTELIGENCIA ARTIFICIAL
%% ——————————————————————————————————————————————————————

\section{Motor de Inteligencia Artificial — Agente Oracle}

El Agente Oracle es la demostración técnica más directa de la profundidad de FABRIC.
Antes de hablar con cualquier consultor, el CTO puede interrogarlo para verificar
si el conocimiento de la firma es real o superficial.

\subsection{Estado de implementación}

\begin{table}[h]
\centering
\begin{tabularx}{\linewidth}{X >{\small\bfseries}l}
\toprule
Componente & Estado \\
\midrule
UI de la terminal (3 escenarios, typewriter, chips) & \color{accent}Implementado \\
Escenario 1: Rescate post go-live & \color{accent2}Respuesta hardcodeada \\
Escenario 2: Migración SAP/EBS/JDE/PeopleSoft & \color{accent2}Respuesta hardcodeada \\
Escenario 3: Greenfield Oracle Fusion & \color{accent2}Respuesta hardcodeada \\
Disclaimer obligatorio al final de respuesta & \color{accent}Implementado \\
Anthropic API (Claude Sonnet 4.6) & \color{danger}Pendiente \\
RAG sobre base de conocimiento FABRIC & \color{danger}Pendiente \\
Kill Switch por jailbreak & \color{danger}Pendiente \\
High-Intent Routing a Office Hours & \color{danger}Pendiente \\
Fallback GPT-4o → Grok & \color{danger}Pendiente \\
Captura de email corporativo opcional & \color{danger}Pendiente \\
Log de conversaciones en MongoDB & \color{danger}Pendiente \\
\bottomrule
\end{tabularx}
\end{table}

\begin{pendiente}
La implementación actual en \ruta{src/pages/public/chat/chatIa.tsx}
no realiza ninguna llamada HTTP. Los 3 escenarios retornan texto prefijado.
No hay integración con la Anthropic API. Esto es \textbf{V0.5 (demo funcional)}
— la arquitectura de la pantalla está lista, falta el motor real.
\end{pendiente}

\subsection{Especificación del motor real}

\begin{table}[h]
\centering
\begin{tabularx}{\linewidth}{>{\ttfamily\color{accent}\small}l X}
\toprule
Parámetro & Especificación \\
\midrule
Modelo primario    & Claude Sonnet 4.6 — Anthropic API \\
Fallback 1         & OpenAI GPT-4o — si Claude supera 3s o devuelve error \\
Fallback 2         & Grok — si GPT-4o falla \\
Streaming          & Respuesta visible carácter a carácter vía EventSource \\
Temperatura        & 0.3 — respuestas precisas y consistentes \\
Scope              & Solo Oracle: Fusion Cloud, EBS, JDE, PeopleSoft, OCI, OIC \\
Kill Switch        & Aborta respuestas fuera de scope o jailbreak \\
High-Intent        & Detecta urgencia y activa CTA de Office Hours \\
RAG                & Búsqueda semántica en embeddings Oracle + material interno FABRIC \\
\bottomrule
\end{tabularx}
\end{table}

\subsection{System Prompt oficial V1.0}

Definido por Julio Álvarez en el Brief 2. Se copia directamente al endpoint de API
sin modificaciones. Ver \ruta{docs/Brief2.md} §7 para el texto completo.

Los puntos clave del system prompt:
\begin{itemize}[leftmargin=*]
  \item Identidad: FABRIC AI Diagnostic, especializado en Oracle
  \item Scope estricto: rechaza educadamente temas fuera de Oracle
  \item Tono: profesional, técnico, directo. Sin emojis ni exclamaciones
  \item Estructura de respuesta: (1) respuesta directa · (2) detalle técnico ·
        (3) riesgos · (4) CTA a FABRIC
  \item 5 reglas críticas: no inventar · no prometer contratos · no info confidencial ·
        redirigir a humano · no criticar competencia
  \item Disclaimer obligatorio al final de cada respuesta
  \item Captura de correo corporativo: sugerir una vez, no insistir
\end{itemize}

\subsection{Implementación del endpoint}

El scaffolding del backend en \ruta{backend/src/} está listo para recibir
la implementación del Agente Oracle:

\begin{lstlisting}[caption={backend/src/routes/ai.ts — a implementar}]
// POST /api/ai
// Body: { message: string, sessionId: string }
// Response: text/event-stream (streaming)
import { Anthropic } from "@anthropic-ai/sdk";

const client = new Anthropic({ apiKey: process.env.ANTHROPIC_API_KEY });

router.post("/", async (req, res) => {
  const { message, sessionId } = req.body;
  // 1. Scope filter
  // 2. RAG lookup
  // 3. Claude Sonnet 4.6 streaming
  // 4. High-Intent detection
  // 5. Log to MongoDB
});
\end{lstlisting}

\subsection{Contenido RAG pendiente}

\begin{table}[h]
\centering
\begin{tabularx}{\linewidth}{X >{\small\bfseries\color{danger}}l}
\toprule
Categoría & Estado \\
\midrule
Documentación Oracle Fusion Cloud pública (Financials, SCM, HCM) & Pendiente \\
Casos FABRIC: APE Plazas y Aplazo con detalle técnico & Pendiente · autorización cliente \\
Doctrina FABRIC V4.0 con contexto operativo & Disponible en Brief2.md \\
Papers 01, 02 y 03 — outlines completos & Disponible en Brief2.md \\
Patrones de fracaso Oracle en LATAM & Pendiente · material interno FABRIC \\
6 FSOs con descripción técnica & Pendiente · confirmación Julio \\
System Prompt V1.0 con ejemplos de respuestas & Listo — Brief2.md §7 \\
\bottomrule
\end{tabularx}
\end{table}

\newpage

%% ——————————————————————————————————————————————————————
%%  §8  CENTRO DE MANDO ADMINISTRATIVO
%% ——————————————————————————————————————————————————————

\section{Centro de mando administrativo}

Julio y el equipo de FABRIC pueden operar el sitio — gestionar leads, ajustar
métricas, controlar la capacidad — sin modificar código ni hacer deploys.

\begin{notatec}[Corrección v6.0 — stack del admin]
El documento v5.0 describía el admin dentro de \texttt{app/admin/} de Next.js.
La implementación real usa React + React Router v7 en \ruta{src/pages/admin/}.
La autenticación con Clerk y la estructura funcional son idénticas a lo especificado.
\end{notatec}

\subsection{Acceso y autenticación}

\begin{table}[h]
\centering
\begin{tabularx}{\linewidth}{>{\ttfamily\color{accent}\small}l X}
\toprule
Parámetro & Especificación \\
\midrule
Ruta de acceso & /admin (protegida por Clerk en React Router) \\
Proveedor de auth & Clerk 5.x — sesiones, rotación de tokens, roles \\
Restricción de dominio & Solo @fabricsoft.com.mx \pendientetag \\
Doble factor & Obligatorio para todos \pendientetag \\
Protección de rutas & \texttt{<SignedIn>} wrapper en cada página admin \\
Roles & Super Admin (Julio): acceso total · Admin: lectura y gestión \\
Audit trail & Cada acción registrada en FabricContext \pendientetag{} (MongoDB pendiente) \\
\bottomrule
\end{tabularx}
\end{table}

\subsection{Funcionalidades implementadas}

\subsubsection{Dashboard (\ruta{AdminDashboard.tsx})}

\begin{itemize}[leftmargin=*]
  \item Fichas de estado: leads nuevos hoy, pendientes de revisión
  \item N/12 proyectos activos
  \item Slots de Office Hours disponibles
  \item Últimos 5 leads con acciones rápidas
\end{itemize}

\subsubsection{Gestión de leads (\ruta{AdminLeads.tsx})}

\begin{itemize}[leftmargin=*]
  \item Vista tabular: nombre, empresa, cargo, industria, status, fecha
  \item Filtros por tipo de iniciativa, industria, estado y rango de fechas
  \item Acciones: Aprobar · Rechazar · Ver detalle
  \item \pendientetag{} HubSpot: aprobar crea contacto calificado en CRM
  \item \pendientetag{} Resend: rechazar envía email educado automático
\end{itemize}

\subsubsection{Editor de métricas (\ruta{AdminMetricas.tsx})}

\begin{itemize}[leftmargin=*]
  \item Rescue Counter: rescates, horas ahorradas, reportes eliminados, cierres
  \item Wait List: proyectos activos, lista de espera, próxima ventana
  \item \pendientetag{} Propagación ISR: sin Sanity/MongoDB, los cambios son solo en sesión
\end{itemize}

\subsubsection{Panel de capacidad (\ruta{AdminCapacidad.tsx})}

\begin{itemize}[leftmargin=*]
  \item Variable de capacidad máxima (default: 12 proyectos)
  \item Toggle Sold Out manual
  \item Calendario de admisión Q1-Q4 con estado y fecha límite
  \item Contador regresivo configurable del plazo de aplicación
\end{itemize}

\subsubsection{Office Hours (\ruta{AdminOfficeHours.tsx})}

\begin{itemize}[leftmargin=*]
  \item Tabla de slots: crear, editar, eliminar
  \item \pendientetag{} URL Calendly por slot
  \item \pendientetag{} Webhook Calendly → MongoDB
\end{itemize}

\subsubsection{Logs (\ruta{AdminLogs.tsx})}

\begin{itemize}[leftmargin=*]
  \item Audit trail completo de acciones del admin
  \item \pendientetag{} Append-only en MongoDB (actualmente en memoria)
\end{itemize}

\newpage

%% ——————————————————————————————————————————————————————
%%  §9  INTEGRACIONES, DATOS Y FLUJOS
%% ——————————————————————————————————————————————————————

\section{Integraciones, datos y flujos de captura}

\subsection{Estado de servicios externos}

\begin{table}[h]
\centering
\begin{tabularx}{\linewidth}{>{\bfseries}l X >{\small\bfseries}l}
\toprule
Servicio & Qué hace & Estado \\
\midrule
Clerk 5.x &
  Autenticación del panel admin. Sesiones, roles, middleware &
  \color{accent}Implementado \\
Anthropic API &
  Claude Sonnet 4.6 — modelo primario del Agente Oracle &
  \color{danger}Pendiente \\
OpenAI API &
  GPT-4o — fallback 1 + embeddings para RAG &
  \color{danger}Pendiente \\
MongoDB Atlas &
  Leads, diagnósticos, conversaciones IA, bookings, audit logs &
  \color{danger}Pendiente · en memoria \\
Sanity CMS &
  Casos, papers, FSOs, métricas, slots editables por Julio &
  \color{danger}Pendiente \\
Resend &
  Email transaccional: confirmaciones, PDFs, notificaciones &
  \color{danger}Pendiente \\
HubSpot &
  CRM: lead creado en cada submit con tags de origen &
  \color{danger}Pendiente \\
Calendly &
  Booking de Office Hours con webhook a MongoDB &
  \color{danger}Pendiente \\
Slack Webhook &
  Notificación interna por cada lead nuevo con score &
  \color{danger}Pendiente \\
Cloudflare Turnstile &
  CAPTCHA invisible en todos los formularios &
  \color{danger}Pendiente \\
Plausible Analytics &
  Analítica privacy-first con eventos custom &
  \color{danger}Pendiente \\
Sentry &
  Monitoreo de errores en producción &
  \color{danger}Pendiente \\
\bottomrule
\end{tabularx}
\end{table}

\subsection{Modelo de datos — MongoDB Atlas (especificación)}

MongoDB fue elegido por su flexibilidad para datos semiestructurados
(conversaciones del agente IA, respuestas variables del Rescue Assessment).

\begin{table}[h]
\centering
\begin{tabularx}{\linewidth}{>{\ttfamily\color{accent2}\small}l X}
\toprule
Colección & Propósito y campos clave \\
\midrule
leads &
  contact, company, qualification, source, score, status, tags \\
diagnosticRuns &
  leadId, answers, painSignals, score, recommendation, submittedAt \\
tcoComparisons &
  leadId, currentStack, annualCost, savingsEstimate, breakevenMonths \\
aiConversations &
  leadId, sessionId, messages, scopeStatus, provider, highIntent \\
paperRequests &
  leadId, paperSlug, email, company, consent, deliveredAt \\
officeHourBookings &
  leadId, calendlyEventId, slot, status, qualificationSnapshot \\
capacityState &
  activeProjects, maxProjects, waitListCount, forcedMode \\
publicMetrics &
  rescueCount, hoursSaved, npsScore, updatedAt \\
adminAuditLogs &
  adminUserId, action, targetCollection, before, after, createdAt (append-only) \\
\bottomrule
\end{tabularx}
\end{table}

\subsection{Flujos de captura — estado real}

\subsubsection{Flujo 1 — Rescue Assessment (/diagnostico) \parcialtag}

UI existe en \ruta{s07b-rescue-assessment.tsx}. Preguntas sobre síntomas Oracle fallido,
score en navegador, output inmediato. \textbf{Submit → backend pendiente}. Sin Zod, sin
MongoDB, sin HubSpot, sin Resend.

\subsubsection{Flujo 2 — ERP TCO Comparator \implementadotag\ (UI)}

UI funcional en \ruta{s03-tco-calculator.tsx}: modal con 8 campos, cálculo TCO en
tiempo real. \textbf{Gating por correo corporativo y entrega de PDF: pendiente}.

\subsubsection{Flujo 3 — Cloud Cost Comparator \parcialtag}

UI en \ruta{s02-optimizador.tsx}. Comparador OCI vs. AWS/GCP/Azure.
\textbf{Análisis personalizado con datos reales: pendiente}.

\subsubsection{Flujo 4 — Wait List (/aplicar) \pendientetag}

Sección S15 muestra la capacidad y el Wait List. \textbf{La ruta
\ruta{/aplicar} con el formulario de 5 campos no existe todavía.}
Este es el CTA final de todo el sitio.

\subsubsection{Flujo 5 — Office Hours \parcialtag}

UI de calendario en \ruta{s11-office-hours.tsx}. \textbf{Sin Calendly real ni gate
de admisión funcional}.

\subsubsection{Flujo 6 — Papers Técnicos \parcialtag}

UI de 3 papers en \ruta{s14-investigacion.tsx}. \textbf{Sin PDFs reales, sin
gating por correo, sin Resend}. Botón dirige a \texttt{\#aplicar}.

\subsection{Arquitectura del backend (scaffolding)}

Se creó la carpeta \ruta{backend/} con estructura Express + TypeScript lista
para implementar:

\begin{lstlisting}[caption={backend/ — estructura actual}]
backend/
|-- src/
|   |-- index.ts           # Express app, CORS, rutas
|   |-- routes/
|   |   |-- leads.ts       # GET /api/leads, POST, PATCH /:id
|   |   |-- metricas.ts    # GET /api/metricas (publico), PATCH /:id
|   |   `-- officeHours.ts # GET slots, POST /:id/reservar, PATCH /:id
|   |-- middleware/
|   |   `-- auth.ts        # requireAuth - API key (Clerk backend pendiente)
|   `-- models/
|       `-- types.ts       # Lead, OfficeHoursSlot, MetricaPublica (sync con frontend)
|-- package.json           # Express, cors, dotenv, tsx
|-- tsconfig.json
`-- .env.example           # PORT, FRONTEND_URL, API_KEY
\end{lstlisting}

\newpage

%% ——————————————————————————————————————————————————————
%%  §10  INFRAESTRUCTURA Y DESPLIEGUE
%% ——————————————————————————————————————————————————————

\section{Infraestructura y despliegue}

\subsection{Entorno de desarrollo}

\begin{table}[h]
\centering
\begin{tabularx}{\linewidth}{>{\ttfamily\color{accent2}}l X}
\toprule
Comando & Descripción \\
\midrule
\texttt{pnpm dev} & Frontend Vite en \texttt{http://localhost:5173} \\
\texttt{cd backend \&\& pnpm dev} & Backend Express en \texttt{http://localhost:3001} \\
\texttt{cd frontend \&\& pnpm dev} & Next.js (Edilberto) en \texttt{http://localhost:3000} \\
\texttt{npx tsc --noEmit} & TypeScript check — sin output = sin errores \\
\bottomrule
\end{tabularx}
\end{table}

\subsection{Repositorio}

\begin{itemize}[leftmargin=*]
  \item Repo: \url{https://github.com/Tiboryeah/FabricSoftPage}
  \item Rama principal: \texttt{main}
  \item Rama del compañero: \texttt{feature/fabric\_edilberto}
  \item CI: no configurado todavía
\end{itemize}

\subsection{Pendientes de infraestructura}

\begin{itemize}[leftmargin=*]
  \item \pendientetag{} Hosting para el frontend Vite (Vercel / Cloudflare Pages)
  \item \pendientetag{} Hosting para el backend Express (Railway / Render / Fly.io)
  \item \pendientetag{} MongoDB Atlas — cuenta y colecciones
  \item \pendientetag{} Variables de entorno para producción
  \item \pendientetag{} CI/CD con GitHub Actions
  \item \pendientetag{} Dominio \texttt{fabricsoft.com.mx} apuntado al hosting
  \item \pendientetag{} SSL, Cloudflare DDoS, headers HTTP de seguridad
\end{itemize}

\subsection{Estrategia de performance — objetivos Lighthouse}

\begin{table}[h]
\centering
\begin{tabularx}{\linewidth}{>{\bfseries}l >{\ttfamily\color{accent}}l X}
\toprule
Métrica & Objetivo & Estado \\
\midrule
Performance & 95+ & No medido en producción todavía \\
LCP         & $<$ 1.5s & Pendiente de deploy \\
CLS         & $<$ 0.05 & Pendiente de deploy \\
SEO         & 100 & Pendiente — meta tags manuales \\
Accesibilidad & 95+ & Pendiente de auditoría \\
\bottomrule
\end{tabularx}
\end{table}

\subsection{Seguridad}

\begin{itemize}[leftmargin=*]
  \item Clerk protege todas las rutas \texttt{/admin/*}
  \item API Key en \ruta{backend/.env} para endpoints del admin
  \item \pendientetag{} Rate limiting (3 intentos/hora/IP en formularios)
  \item \pendientetag{} Cloudflare Turnstile en formularios públicos
  \item \pendientetag{} Headers HTTP: CSP, HSTS, X-Frame-Options
  \item \pendientetag{} Validación Zod server-side para todos los endpoints
\end{itemize}

\newpage

%% ——————————————————————————————————————————————————————
%%  RESUMEN EJECUTIVO DE COBERTURA
%% ——————————————————————————————————————————————————————

\section*{Resumen ejecutivo de cobertura del Brief 2}
\addcontentsline{toc}{section}{Resumen ejecutivo de cobertura}

\begin{table}[h]
\centering
\begin{tabularx}{\linewidth}{X >{\bfseries\small}l}
\toprule
Área & Estado \\
\midrule
Posicionamiento y doctrina (§1, §2) & \color{accent}100\% implementado \\
Sistema de diseño — tokens, paleta, microinteracciones (§4) & \color{accent}95\% implementado \\
19 secciones del home (§3) & \color{accent2}80\% implementado · 20\% UI sin backend \\
Rutas activas & \color{accent2}10 de 15 rutas activas \\
Motor IA — UI (§7) & \color{accent}UI completa \\
Motor IA — API real Claude Sonnet 4.6 (§7) & \color{danger}Pendiente \\
Panel admin funcional (§8) & \color{accent}Implementado · persistencia pendiente \\
Integraciones externas (§9) & \color{danger}Solo Clerk — resto pendiente \\
PDFs de papers + gating (§9 Flujo 6) & \color{danger}Pendiente \\
Formulario /aplicar (§9 Flujo 4) & \color{danger}Pendiente — crítico \\
Backend Express scaffolding (§10) & \color{accent}Estructura lista \\
Persistencia MongoDB (§10) & \color{danger}Pendiente \\
Deploy a producción (§10) & \color{danger}Pendiente \\
\bottomrule
\end{tabularx}
\caption{Estado de cobertura · Mayo 2026 · post-demo}
\end{table}

\subsection*{Pendientes críticos para V1 en producción}

\begin{enumerate}[leftmargin=*, label=\textbf{\arabic*.}]
  \item \textbf{Ruta \ruta{/aplicar} con formulario de 5 campos} —
        Es el CTA final de todo el sitio. Sin esto el visitante no tiene a dónde ir.
  \item \textbf{Agente Oracle con Anthropic API real} —
        El endpoint en \ruta{backend/src/routes/} está scaffoldeado.
        Requiere API key de Anthropic y \texttt{VITE\_API\_URL} apuntando al backend.
  \item \textbf{MongoDB + persistencia de leads} —
        El admin panel funciona pero los datos se pierden al recargar.
  \item \textbf{PDFs de los 3 papers + gating por correo corporativo} —
        Los papers son la credencial académica del sitio; sin descarga real,
        la sección es un placeholder.
  \item \textbf{Deploy a producción en \texttt{fabricsoft.com.mx}} —
        Verificar Lighthouse 95+ en producción real.
\end{enumerate}

\newpage

%% ——————————————————————————————————————————————————————
%%  ANEXO A — BRIEF FOTOGRÁFICO
%% ——————————————————————————————————————————————————————

\section*{Anexo A — Brief fotográfico Julio Álvarez}
\addcontentsline{toc}{section}{Anexo A · Brief fotográfico}

Esta es la única fotografía del sitio. Todo el peso visual de la sección S15
descansa en ella.

\subsection*{Concepto}

\textit{``El ingeniero que asume el riesgo''}

Referencias: páginas de partners en Bain \& Company, Sequoia Capital,
Founders Fund, Stripe leadership.

Mirada directa a cámara con autoridad tranquila. Postura natural.
Expresión seria pero accesible (10--20\% de sonrisa con ojos, no con boca).
Sin pose corporativa de LinkedIn.

\subsection*{Especificaciones técnicas para el fotógrafo}

\begin{table}[h]
\centering
\begin{tabularx}{\linewidth}{>{\bfseries}l X}
\toprule
Parámetro & Detalle \\
\midrule
Tipo & Editorial corporativo premium · B\&W o color muy desaturado \\
Formatos & Retrato medio (pecho) · retrato 3/4 (cintura) · manos opcional \\
Orientaciones & Vertical (mobile) · cuadrada (desktop) · horizontal amplia \\
Fondo & Estudio: gris carbón \texttt{\#2B2B2B} o negro puro. Sin texturas. \\
Iluminación & Lateral suave (Rembrandt o split). Sombras controladas. \\
Vestimenta A & Camisa sólida oscura. Sin corbata. Sin accesorios. \\
Vestimenta B & Saco negro estructurado. Camisa oscura debajo. Sin corbata. \\
Post-producción & Editorial · textura natural · sin filtros · B\&W + color \\
Entregables & 5-8 seleccionadas · edición de 3 · RAW + JPG + PNG · derechos perpetuos \\
\bottomrule
\end{tabularx}
\end{table}

\textbf{Prohibiciones absolutas:} camisa blanca con corbata · polos · trajes con
patrones · camisas con logo · lentes de sol · joyería visible.

\newpage

%% ——————————————————————————————————————————————————————
%%  ANEXO B — PAPER APE PLAZAS
%% ——————————————————————————————————————————————————————

\section*{Anexo B — Paper APE Plazas (borrador para producción editorial)}
\addcontentsline{toc}{section}{Anexo B · Paper APE Plazas}

\begin{pendiente}
Publicación pendiente de: (1) autorización formal de APE Plazas,
(2) revisión legal del NDA mutuo. El diseño editorial final se produce en
Figma o InDesign con el mismo sistema de diseño del sitio.
\end{pendiente}

\vspace{1em}

\begin{center}
{\large\bfseries FABRIC RESEARCH · CASO 01}\\[0.4em]
\textit{Implementación Oracle Fusion Cloud en Operadora de Centros Comerciales}\\[0.2em]
{\small Entrega en primer ciclo crítico · Abril 2026 · Publicado con autorización
del cliente bajo NDA mutuo}
\end{center}

\vspace{1em}
\hrule
\vspace{1em}

\subsubsection*{Resumen ejecutivo}

En abril de 2026, FABRIC completó la implementación de Oracle Fusion Cloud para
una operadora de centros comerciales en México. El proyecto no fue considerado
entregado en el go-live del 6 de abril. Fue considerado entregado el 30 de abril,
al cierre del primer ciclo contable ejecutado en producción con acompañamiento
directo de FABRIC. Ese es el estándar de la firma.

\subsubsection*{Ejecución por fase}

\begin{table}[h]
\centering
\begin{tabularx}{\linewidth}{>{\ttfamily\color{accent}\bfseries}l X}
\toprule
Fase & Qué se hizo \\
\midrule
01 · DIAGNOSE   & Análisis de procesos financieros, riesgos de integración y mapeo de stakeholders \\
02 · ARCHITECT  & Diseño técnico multi-entidad, cuentas contables, integraciones con bancos \\
03 · DEPLOY     & Implementación FABRIC con control de cambios y gobierno semanal. Go-live: 06 abril 2026 \\
04 · STABILIZE  & Soporte presencial y remoto durante abril. Validación de cierre día a día. Reuniones diarias con CFO. \\
05 · OPTIMIZE   & Plan de optimización trimestral acordado. Activo. \\
\bottomrule
\end{tabularx}
\end{table}

\subsubsection*{Resultados verificados}

\begin{table}[h]
\centering
\begin{tabularx}{\linewidth}{X >{\color{accent}\bfseries}l}
\toprule
Hito & Estado \\
\midrule
Go-live en fecha contractual: 06 abril 2026 & Cumplido \\
Configuración multi-entidad en producción   & Operativa \\
Cierre contable de abril 2026 sin incidencias críticas & Ejecutado el 30 de abril \\
Tablero ejecutivo activo para Dirección y CFO & Activo \\
Adopción de usuarios clave documentada & Verificada \\
Acta de transición a soporte & Firmada por todos los stakeholders \\
\bottomrule
\end{tabularx}
\end{table}

\subsubsection*{Lecciones documentadas}

\begin{enumerate}[leftmargin=*]
  \item \textbf{El go-live no es el riesgo. El primer cierre lo es.} El riesgo
        operativo real aparece en los primeros 15--30 días en producción, cuando
        los procesos de negocio se ejecutan bajo carga real por primera vez.
  \item \textbf{Una célula multidisciplinaria acelera la estabilización.}
        La combinación de consultor senior Oracle, especialista en procesos
        financieros y especialista en IA permite resolver incidencias en horas,
        no en días.
  \item \textbf{La transición a soporte debe ser un hito formal.}
        Un acta firmada por todos los stakeholders elimina cualquier ambigüedad
        sobre cuándo termina el proyecto.
\end{enumerate}

\vspace{2em}

\begin{center}
\small\ttfamily\color{textsec}
FABRIC Oracle Critical Engineering · Análisis y Diseño V6.0 · Mayo 2026 \\
Equipo A · Uso interno · Para revisión de Julio Álvarez
\end{center}

\end{document}
